\subsection{The topology defined by a filtration}

Let \(\mathbf{G}\) be a group filtered by a family \((\mathbf{G}_n)_{n\in \mathbf{Z}}\) of normal subgroups of \(\mathbf{G}\) . There exists a unique topology on \(\mathbf{G}\) which is compatible with the group structure and for which the \(G_{n}\) constitute a fundamental system of neighbourhoods of the identity element e of G (General Topology, Chapter III, § 1, no. 2, Example); it is called the topology on G defined by the filtration ( \(G_{n}\) ). When we use topological notions concerning a filtered group, we shall mean, unless otherwise stated, with the topology defined by the filtration. Note that the \(G_{n}\) , being subgroups of G, are both open and closed (General Topology, Chapter III, § 2, no. 1, Corollary to Proposition 4).

As each \(G_{n}\) is normal in G, the entourages of the left and right uniformities on G coincide; we deduce that G admits a Hausdorff completion group \(\hat{G}\) (General Topology, Chapter III, § 3, no. 4, Theorem 1 and no. 1, Proposition 2).

For every subset M of G, the closure of M in G is equal to

\[
\bigcap_ {n \in \mathbf {Z}} (\mathbf {M}. \mathbf {G} _ {n}) = \bigcap_ {n \in \mathbf {Z}} (\mathbf {G} _ {n}. \mathbf {M})
\]

(General Topology, Chapter III, § 3, no. 1, formula (1)); in particular \(\bigcap_{n\in\mathbf{Z}}\mathrm{G}_{n}\) is the closure of \(\{e\}\) ; thus it is seen that for the topology on \(\mathbf{G}\) to be Hausdorff it is necessary and sufficient that the filtration \((\mathrm{G}_{n})\) be separated. For the topology on \(\mathbf{G}\) to be discrete, it is necessary and sufficient that there exist \(n\in\mathbf{Z}\) such that \(\mathrm{G}_{n}=\{e\}\) (in which case \(\mathrm{G}_{m}=\{e\}\) for \(m\geqslant n\) ); then the filtration \((\mathrm{G}_{n})\) is called discrete.

Since the Hausdorff group associated with G is \(H = \frac{G}{\left(\bigcap_{n \in \mathbf{Z}} G_n\right)}\) , the associated graded groups \(\text{gr}(G)\) and \(\text{gr}(H)\) (if H is given the quotient filtration) are canonically identified.

Not let \(G'\) be another filtered group and \(u: G \to G'\) a homomorphism compatible with the filtrations; the definition of the topologies on G and \(G'\) shows immediately that u is continuous (*). If H is a subgroup (resp. normal subgroup) of G, the topology induced on H by that on G (resp. the quotient topology with respect to H of that on G) is the topology on H (resp. G/H) defined by the filtration induced by that on G (resp. quotient topology of that on G). The product topology of those on G and \(G'\) is the topology defined by the product of filtrations on G and \(G'\) .

Let \(v\) be the order function (no. 2) on \(G\) . The hypothesis on the \(G_n\) implies that \(v(xyx^{-1}) = v(y)\) and hence \(v(xy^{-1}) = v(yx^{-1}) = v(x^{-1}y) = v(y^{-1}x)\) for all \(x, y\) in \(G\) . Let \(\rho\) be a real number such that \(0 < \rho < 1\) (for example take \(\rho = 1 / e)\) and let \(d(x,y) = \rho^{y(xy - 1)}\) for all \(x,y\) in G. Then \(d(x,x) = 0, d(x,y) = d(y,x)\) and inequality (5') of no. 2 gives

\[
d (x, y) \leqslant \sup (d (x, z), d (y, z))\tag{15}
\]

for all x, y, z in G, which implies the triangle inequality

\[
d (x, y) \leqslant d (x, z) + d (y, z).
\]

Thus \(d\) is a pseudometric on G which is invariant under left and right translations and \(\mathbf{G}_n\) is the set of \(x \in \mathbf{G}\) such that \(d(e, x) \leqslant \rho^n\) ; the uniform structure defined by \(d\) is then the uniform structure on the topological group G. If G is Hausdorff, G is a zero-dimensional metrizable topological space (General Topology, Chapter IX, § 6, no. 4); \(d\) is a distance on G if also the filtration \((\mathbf{G}_n)\) is exhaustive.

Given a topological ring A, recall that a left topological A-module is an A-module E with a topology compatible with its additive group structure and such that the mapping \((a, x) \mapsto ax\) from A × E to E is continuous (General Topology, Chapter III, § 6, no. 6).

PROPOSITION 3. Let A be a filtered ring, \((\mathbf{A}_n)\) its filtration, B the subring \(\bigcup_{n\in \mathbf{Z}}\mathbf{A}_n\) of A, E a filtered B-module, \((\mathrm{E}_n)\) its filtration and F the sub-B-module \(\bigcup_{n\in \mathbf{Z}}\mathrm{E}_n\) of E. Then the mapping \((a,x)\mapsto ax\) from \(\mathbf{B}\times \mathbf{F}\) to F is continuous.

Let \(a_0 \in \mathbf{B}\) , \(x_0 \in \mathbf{F}\) ; there exists by hypothesis integers \(r, s\) such that \(a_0 \in \mathbf{A}_r\) and \(x_0 \in \mathbf{E}_s\) . The relation

\[
a x - a _ {0} x _ {0} = (a - a _ {0}) x _ {0} + a _ {0} (x - x _ {0}) + (a - a _ {0}) (x - x _ {0})
\]

shows that if \(a - a_0 \in \mathbf{A}_l\) and \(x - x_0 \in \mathrm{E}_f\) , \(ax - a_0x_0\) belongs to

\[
\mathbf {E} _ {i + s} + \mathbf {E} _ {j + r} + \mathbf {E} _ {i + j}.
\]

Then, given an integer \(n\) , \(ax - a_0x_0 \in \mathrm{E}_n\) provided \(i \geqslant n - s\) , \(j \geqslant n - r\) and \(i + j \geqslant n\) , that is so long as \(i\) and \(j\) are sufficiently large.

COROLLARY. The ring B is a topological ring and the B-module F a topological B-module.

The first assertion is obtained by applying Proposition 3 to \(\mathbf{F} = \mathbf{B}_s\) .

It is seen in particular that a filtered ring A whose filtration is exhaustive is a topological ring; if this is so every filtered A-module whose filtration is exhaustive is a topological A-module.

PROPOSITION 4. Let A be a commutative ring filtered by an exhaustive filtration \((\mathbf{A}_n)\) and \(\mathfrak{p}\) an ideal of A. Suppose that the ideal \(\operatorname{gr}(\mathfrak{p}) = \bigoplus_{n \in \mathbf{Z}} (\mathfrak{p} \cap \mathbf{A}_n) / (\mathfrak{p} \cap \mathbf{A}_{n+1})\) of the ring \(\operatorname{gr}(\mathbf{A})\) is prime. Then the closure of \(\mathfrak{p}\) in A is a prime ideal.

We know that \(\operatorname{gr}(\mathbf{A} / \mathfrak{p})\) is isomorphic to \(\operatorname{gr}(\mathbf{A}) / \operatorname{gr}(\mathfrak{p})\) (no. 4, Proposition 2) and hence an integral domain; we conclude that \(\mathbf{A} / \bigcap_{n\in \mathbf{Z}}(\mathfrak{p} + \mathbf{A}_n)\) is an integral domain (no. 3, Corollary to Proposition 1). Then the closure \(\bigcap_{n\in \mathbf{Z}}(\mathfrak{p} + \mathbf{A}_n)\) of \(\mathfrak{p}\) is a prime ideal.

Let A be a ring and m a two-sided ideal of A; the topology defined on A by the m-adic filtration (no. 1, Example 3) is called the m-adic topology; as the m-adic filtration is exhaustive, A is a topological ring with this topology (Corollary to Proposition 3). Similarly, for every A-module E, the topology defined by the m-adic filtration is called the m-adic topology on E; E is a topological A-module under this topology.

Let \(m'\) be another two-sided ideal of A; for the \(m'\) -adic topology on A to be finer than the \(m\) -adic topology, it is necessary and sufficient that there exist an integer n \textgreater{} 0 such that \(m'^n \subset m\) ; the condition is necessary and, if it is fulfilled, \(m'^{hn} \subset m^h\) for all h \textgreater{} 0 and hence the condition is sufficient. If A is a commutative Noetherian ring, it amounts to the same to say that \(V(m) \subset V(m')\) in the prime spectrum of A (Chapter II, § 4, no. 3, Corollary 2 to Proposition 11 and § 2, no. 6, Proposition 15).

